\documentclass[letterpaper]{article}
\usepackage[letterpaper,margin=0.5in]{geometry}
\usepackage[T1]{fontenc}
\usepackage[scaled=1.0]{tgheros}
\renewcommand{\familydefault}{\sfdefault}
\usepackage{ragged2e}
\usepackage{titlesec}
\setlength{\parindent}{0pt}
\setlength{\parskip}{0.35em}
\titleformat{\section}{\bfseries\fontsize{11pt}{12pt}\selectfont}{}{0pt}{}
\titlespacing*{\section}{0pt}{0pt}{0.25em}

\begin{document}
\fontsize{11pt}{12pt}\selectfont
\RaggedRight

\section*{SPECIFIC AIMS}

% Refs for burden and clinical framing: PMID 29097296, PMID 31075787, PMID 37283271.
% Refs for ECGi / inverse limitations: PMID 34777827, PMID 15649241.
Ventricular arrhythmias and conduction disorders drive sudden cardiac death, progressive heart failure, recurrent hospitalizations, implantable cardioverter-defibrillator therapies, and repeat invasive procedures. Yet key decisions about ventricular tachycardia ablation, pacing strategy, and evaluation of inherited arrhythmia syndromes are still made from disconnected snapshots provided by surface ECG, cardiac imaging, and invasive electrophysiology studies. These tests rarely explain, in an individual patient, how activation propagates through three-dimensional ventricular anatomy, which substrate abnormalities produce the observed phenotype, or how activation is likely to change with pacing or other interventions. Electrocardiographic imaging (ECGi) partially addresses this gap, but it reconstructs electrical activity only on the cardiac surface from an ill-posed inverse problem and does not robustly identify internal substrate or distinguish alternative mechanisms that can produce similar body-surface signals.

% Refs for translational modeling: PMID 33303478, PMID 38323181.
% Refs for uncertainty-aware calibration: PMID 32448065.
The long-term goal of this work is to enable non-invasive, mechanism-based precision electrophysiology. The objective of this application is to develop and prospectively test a patient-specific fast forward ECG (ffECG) platform that integrates cardiac imaging, body-surface potentials, and mechanistic modeling to infer ventricular activation and substrate non-invasively. Our central hypothesis is that forward-calibrated, uncertainty-aware patient-specific models will recover clinically meaningful activation patterns and substrate abnormalities more faithfully than inverse ECGi or conventional ECG and imaging metrics, and that these computational phenotypes will improve localization, mechanistic classification, and prediction of therapy response. This hypothesis is supported by our preliminary work in patient-specific forward modeling and calibration, including a research ffECG workflow, curated retrospective imaging and body-surface potential datasets, MRI-compatible ECGi acquisition, and emerging GPU-native simulation and emulation methods. The rationale for the proposed work is that a fast, automated, and uncertainty-aware forward-calibrated framework would provide clinically actionable information that is not available from current non-invasive assessment.

We will test our central hypothesis through three complementary aims that develop one common platform and evaluate it across three clinically important electrophysiology testbeds.

\textbf{Aim 1. Develop and validate an automated patient-specific ffECG pipeline for rapid construction of ventricular and torso models from cardiac imaging and body-surface potentials.} We will create a scalable image-to-model workflow that integrates anatomical segmentation, scar and fibrosis labeling, meshing, fiber assignment, and body-surface potential preprocessing into simulation-ready patient-specific models. We will implement GPU-native forward simulation and harmonization tools that are robust to heterogeneous clinical datasets, and benchmark each component against our current reference pipeline in retrospective cohorts spanning healthy volunteers, cardiomyopathy, and ventricular arrhythmia populations. The deliverable for this aim is an automated and validated model-construction and simulation pipeline suitable for translational electrophysiology studies.

% Refs for fibrosis/substrate and translational simulations: PMID 25368538, PMID 33303478.
% Refs for Brugada mechanism/substrate: PMID 31993492, PMID 33421051, PMID 36542434.
\textbf{Aim 2. Infer and validate mechanistic computational phenotypes of ventricular activation and substrate using Bayesian calibration of patient-specific ffECG models.} We will develop an inference framework that combines emulator-accelerated simulation, dimensionality reduction of body-surface potentials, sensitivity analysis, and staged Bayesian calibration to estimate activation sequence, initiation site, conduction slowing, repolarization heterogeneity, and substrate-associated parameters from non-invasive recordings. We will apply this framework in retrospective cohorts with ventricular tachycardia, pacing, and Brugada syndrome to define computational phenotypes linked to arrhythmia substrate, paced electrical response, and inherited arrhythmia mechanisms. These phenotypes will be evaluated against invasive mapping, pacing response, and genotype-phenotype structure where appropriate. The deliverable for this aim is a validated set of mechanistically interpretable digital phenotypes that move beyond descriptive ECG features toward biologically meaningful classification.

% Refs for VT ablation: PMID 31075787.
% Refs for pacing/CRT/CSP: PMID 37283271, PMID 35715087, PMID 37767743.
% Refs for Brugada subgrouping / conduction reserve: PMID 33421051, PMID 36542434.
\textbf{Aim 3. Prospectively test whether ffECG-derived computational phenotypes improve localization and prediction of therapy response in clinical electrophysiology populations.} We will prospectively evaluate ffECG in three complementary settings: patients undergoing ventricular tachycardia ablation, patients undergoing CRT or conduction-system pacing assessment, and patients evaluated for Brugada syndrome. In ventricular tachycardia, we will test whether ffECG localizes abnormal activation and substrate relative to invasive electroanatomic mapping. In pacing populations, we will test whether baseline-calibrated patient-specific models predict paced QRS response and identify favorable pacing configurations. In Brugada syndrome, we will test whether inferred conduction and substrate abnormalities distinguish mechanistically relevant subgroups beyond standard clinical classification. Across settings, performance will be compared directly with current ECG and imaging metrics used in practice. The deliverable for this aim is prospective evidence on whether mechanistic computational phenotyping improves clinically relevant localization and response prediction.

The proposed work is significant because it addresses a central limitation in current electrophysiology practice: the absence of a non-invasive method that can explain patient-specific ventricular activation and substrate in mechanistic terms. It is innovative because it replaces purely inverse reconstruction with a forward-calibrated, uncertainty-aware framework that integrates anatomy, substrate, conduction, and electrical recordings within a single patient-specific pipeline. The expected impact is a translational platform for precision electrophysiology that can improve selection for ablation and pacing therapies, refine mechanistic classification in inherited arrhythmia syndromes, and create a rigorous foundation for broader clinical testing of computationally guided arrhythmia care.

\end{document}
